\chapter{Álgebras \texorpdfstring{$C^*$}{C*} y cálculo funcional continuo}\label{cap:v3-13}

El adjunto de un operador no es una operación auxiliar: vincula su estructura algebraica con su norma. La identidad $\norm{a^*a}=\norm a^2$ recoge ese vínculo y tiene consecuencias muy fuertes: determina la norma, permite reconocer el espectro de elementos normales y convierte un álgebra conmutativa en un espacio de funciones continuas.

\section{Involución e identidad \texorpdfstring{$C^*$}{C*}}\label{lec:v3-25}
\fuentelectura{25}{Álgebras C*}

\begin{definicion}
Una involución en un álgebra compleja $A$ es una operación $a\mapsto a^*$ tal que
\[
 (\alpha a+\beta b)^*=\overline\alpha a^*+\overline\beta b^*,\quad
 (ab)^*=b^*a^*,\quad (a^*)^*=a.
\]
Una \emph{álgebra $C^*$} es un álgebra de Banach con involución que satisface $\norm{a^*a}=\norm a^2$ para todo $a$. No se exige identidad multiplicativa.
\end{definicion}
\index{algebra C*@$C^*$-álgebra}\index{involucion@involución}

La segunda propiedad se llama \emph{antimultiplicatividad}: invierte el orden del producto; no afirma $ab=-ba$. Las conjugaciones de los escalares son indispensables en la primera propiedad.

\begin{proposicion}\label{v3:involucion-isometrica}
La involución de un álgebra $C^*$ es isométrica. Si existe identidad, $1^*=1$ y $\norm1=1$ en el álgebra no nula. Además $(a^{-1})^*=(a^*)^{-1}$ para todo invertible.
\end{proposicion}
\begin{proof}
Para $a\ne0$, $\norm a^2=\norm{a^*a}\le\norm{a^*}\norm a$, de donde $\norm a\le\norm{a^*}$. Aplicar esto a $a^*$ da la desigualdad inversa. De $a1=a=1a$ y la sobreyectividad de la involución se obtiene que $1^*$ también es identidad, luego coincide con $1$. La identidad $C^*$ da $\norm1=\norm1^2$, y la norma no es cero. Aplicar la involución a $aa^{-1}=a^{-1}a=1$ prueba la fórmula de inversas.
\end{proof}

\begin{ejemplo}
$\C$, $C(K)$ para $K$ compacto Hausdorff, $C_b(X)$ y $C_0(X)$ para $X$ localmente compacto Hausdorff, y $L^\infty(\mu)$ son álgebras $C^*$ con conjugación puntual. La identidad de norma se reduce a $\norm{|f|^2}_\infty=\norm f_\infty^2$. En un espacio no compacto, $C(X)$ completo, sin restricción de acotación, no lleva en general una norma del supremo finita; ésa es la corrección necesaria al ejemplo original. En $\Lcal(H)$, el adjunto de Hilbert satisface $\norm{T^*T}=\norm T^2$ y proporciona el ejemplo no conmutativo fundamental.

El álgebra de convolución $L^1(\R)$ tiene involución $f^*(t)=\overline{f(-t)}$. Es conjugada lineal e involutiva, y
\[
 (f*g)^*(t)=\int_{\R}\overline{f(s)g(-t-s)}\,ds
 =\int_{\R}\overline{g(-u)}\,\overline{f(u-t)}\,du
 =(g^**f^*)(t)
\]
casi en todas partes, por el cambio $u=t+s$. No es un álgebra $C^*$ con la norma $L^1$. Para $f=\mathbf1_{[0,1]}-\mathbf1_{[1,2]}$, $\norm f_1^2=4$, mientras su autocorrelación $f^**f$ es par y vale $2-3t$ en $0\le t\le1$, $t-2$ en $1\le t\le2$ y cero para $t>2$. Integrando su valor absoluto se obtiene $\norm{f^**f}_1=8/3<4$.
\end{ejemplo}

\begin{definicion}
Un elemento es \emph{normal} si $a^*a=aa^*$, \emph{autoadjunto} si $a=a^*$, \emph{unitario} si $a^*a=aa^*=1$, \emph{positivo} si tiene la forma $b^*b$, y \emph{proyección} si $p=p^*=p^2$.
\end{definicion}

Todo autoadjunto y todo unitario es normal; todo $b^*b$ es autoadjunto y toda proyección es positiva porque $p=p^*p$. La equivalencia entre esta definición algebraica de positividad y la condición espectral $\sigma(a)\subset[0,\infty)$ se demostrará en el capítulo siguiente, sin utilizarla anticipadamente aquí.

\subsection{Norma y espectro de elementos normales}
Por ahora suponemos que $A$ es unitaria. La unitización que construiremos después extenderá las afirmaciones espectrales al caso no unitario.

\begin{proposicion}\label{v3:normal-radio-norma}
Se tiene $\sigma(a^*)=\overline{\sigma(a)}$. Si $a$ es normal, $r(a)=\norm a$. Si $u$ es unitario, $\sigma(u)\subset\mathbb T$; si $a$ es autoadjunto, $\sigma(a)\subset\R$.
\end{proposicion}
\begin{proof}
La fórmula de inversas muestra que $a-\lambda1$ es invertible si y sólo si $a^*-\overline\lambda1$ lo es. Para $a=a^*$, la identidad $C^*$ implica $\norm{a^2}=\norm a^2$, y al iterar se obtiene $\norm{a^{2^n}}=\norm a^{2^n}$. La fórmula de Gelfand da $r(a)=\norm a$. Si $a$ es normal,
\[
 \norm{a^2}^2=\norm{(a^*)^2a^2}
 =\norm{(a^*a)^2}=\norm{a^*a}^2=\norm a^4.
\]
Las potencias de $a$ siguen siendo normales, y la misma iteración demuestra $r(a)=\norm a$.

Si $u$ es unitario, $\norm u^2=\norm1=1$ y también $\norm{u^{-1}}=1$. Para $\lambda\in\sigma(u)$, las cotas espectrales de $u$ y $u^{-1}$ dan $|\lambda|\le1$ y $|\lambda|^{-1}\le1$, de donde $|\lambda|=1$. Si $a=a^*$, la convergencia en norma de la serie exponencial y la continuidad de la involución dan $(e^{ia})^*=e^{-ia}=(e^{ia})^{-1}$. Así $e^{ia}$ es unitario. El mapeo espectral implica $|e^{i\lambda}|=1$ para $\lambda\in\sigma(a)$; como ese módulo es $e^{-\operatorname{Im}\lambda}$, necesariamente $\operatorname{Im}\lambda=0$.
\end{proof}

\begin{teorema}[Fuglede]\label{v3:fuglede}
Si $a$ es normal y $ab=ba$ en un álgebra $C^*$ unitaria, entonces $a^*b=ba^*$.
\end{teorema}
\begin{proof}
La función $F(z)=e^{za^*}be^{-za^*}$ es entera con valores en $A$, como se ve multiplicando sus series localmente uniformes. Como $b$ conmuta con $a$, conmuta con todas sus potencias y exponenciales. Además $a$ conmuta con $a^*$. Por tanto
\[
 F(z)=e^{za^*-\overline z a}\,b\,e^{-za^*+\overline z a}.
\]
El exponente $za^*-\overline z a$ es antiautoadjunto: su adjunto es su negativo. Su exponencial es unitaria, porque la serie y la involución dan adjunto igual a la exponencial opuesta. En consecuencia $\norm{F(z)}\le\norm b$ para todo $z$. Liouville vectorial implica que $F$ es constante. Derivar en cero da $0=F'(0)=a^*b-ba^*$.

Todo elemento también admite una descomposición única $c=h+ik$ con $h,k$ autoadjuntos, dada por $h=(c+c^*)/2$ y $k=(c-c^*)/(2i)$. La normalidad de $c$ equivale a $hk=kh$, como se verifica desarrollando $c^*c-cc^*=2i(hk-kh)$. Estas identidades explican la separación en partes autoadjuntas empleada en la prueba original; la demostración anterior evita ambigüedades entre $z$ y $\overline z$.
\end{proof}

\section{Ejercicios sobre la estructura \texorpdfstring{$C^*$}{C*}}

\begin{ejercicio}\label{ej:v3-l25-1}
Si $a=a^*$, pruebe que $e^{ia}$ es unitario.
\end{ejercicio}
\begin{solucion}
La involución pasa al límite de la serie por ser isométrica, así $(e^{ia})^*=\sum(-ia)^n/n!=e^{-ia}$. Las exponenciales de $ia$ y $-ia$ multiplican a $1$ en ambos órdenes. Por tanto $(e^{ia})^*e^{ia}=e^{ia}(e^{ia})^*=1$. Si $A$ no es unitaria, la exponencial y la afirmación se entienden en su unitización.
\end{solucion}

\begin{ejercicio}\label{ej:v3-l25-2}
Para $a$ normal pruebe: (i) existe $\lambda\in\sigma(a)$ con $|\lambda|=\norm a$; (ii) si $a^n=0$ para algún $n\ge1$, entonces $a=0$.
\end{ejercicio}
\begin{solucion}
(i) El módulo alcanza su máximo en el espectro compacto no vacío, y ese máximo es $r(a)=\norm a$. (ii) La nilpotencia implica $r(a)=\lim_k\norm{a^k}^{1/k}=0$. La normalidad convierte esa igualdad en $\norm a=0$. Sin normalidad la conclusión es falsa: la matriz $\left(\begin{smallmatrix}0&1\\0&0\end{smallmatrix}\right)$ es nilpotente y no nula.
\end{solucion}

\begin{ejercicio}\label{ej:v3-l25-3}
Decimos que $a,b$ son unitariamente equivalentes si $b=uau^*$ para algún unitario $u$. Pruebe: (i) ésta es una relación de equivalencia; (ii) preserva normalidad, autoadjunción y positividad; (iii) $\sigma(uau^*)=\sigma(a)$.
\end{ejercicio}
\begin{solucion}
(i) La identidad prueba reflexividad; si $b=uau^*$, entonces $a=u^*bu$; y si $c=vbv^*$, entonces $c=(vu)a(vu)^*$, con $vu$ unitario. (ii) Para $b=uau^*$ se tiene $b^*=ua^*u^*$, $b^*b=ua^*au^*$ y $bb^*=uaa^*u^*$; las dos primeras propiedades se conservan. Si $a=c^*c$, entonces $b=(cu^*)^*(cu^*)$, y es positivo. (iii) $b-\lambda1=u(a-\lambda1)u^*$; ese elemento es invertible exactamente cuando lo es $a-\lambda1$, porque los factores externos son invertibles.
\end{solucion}

\begin{ejercicio}\label{ej:v3-l25-4}
Pruebe que los unitarios forman un subgrupo cerrado de $G(A)$. Determine si es necesariamente normal.
\end{ejercicio}
\begin{solucion}
El producto de unitarios es unitario y la inversa de $u$ es $u^*$. Si $u_n\to u$ en norma, la continuidad del producto y de la involución da $u^*u=uu^*=1$; el grupo unitario es incluso cerrado en $A$. No tiene por qué ser normal en $G(A)$. En $M_2(\C)$, tome
\[
 u=\begin{pmatrix}0&1\\1&0\end{pmatrix},\qquad
 s=\begin{pmatrix}2&0\\0&1\end{pmatrix}.
\]
Entonces $sus^{-1}=\left(\begin{smallmatrix}0&2\\1/2&0\end{smallmatrix}\right)$ y $(sus^{-1})^*(sus^{-1})=\operatorname{diag}(1/4,4)\ne I$.
\end{solucion}

\begin{ejercicio}\label{ej:v3-l25-5}
Pruebe que $A$ es conmutativa si y sólo si todos sus elementos son normales.
\end{ejercicio}
\begin{solucion}
La implicación directa es inmediata. Recíprocamente, si $h,k$ son autoadjuntos, la normalidad de $h+ik$ y el desarrollo al final de la prueba de Fuglede dan $hk=kh$. Cada elemento es combinación compleja lineal de dos autoadjuntos; distribuir los productos muestra entonces que cualesquiera dos elementos de $A$ conmutan.
\end{solucion}

\begin{ejercicio}\label{ej:v3-l25-6}
Pruebe que una misma álgebra con involución tiene a lo sumo una norma que la convierta en álgebra $C^*$.
\end{ejercicio}
\begin{solucion}
En cualquiera de esas normas, $a^*a$ es autoadjunto, de modo que
\[
 \norm a^2=\norm{a^*a}=r(a^*a).
\]
El miembro derecho sólo depende de la invertibilidad algebraica de $a^*a-\lambda1$, y no de la norma elegida. Por tanto las dos normas coinciden en cada $a$. En el caso no unitario se usa la unitización algebraica $A\oplus\C$, que es la misma para ambas normas; el teorema siguiente garantiza una norma $C^*$ en ella y permite aplicar el mismo argumento sin suponer de antemano que las normas originales sean equivalentes.
\end{solucion}

\begin{ejercicio}\label{ej:v3-l25-7}
Pruebe que un álgebra de Banach con involución es $C^*$ si y sólo si $\norm{a^*a}\ge\norm a^2$ para cada $a$.
\end{ejercicio}
\begin{solucion}
La necesidad es inmediata. Para la suficiencia, si $a\ne0$,
$\norm a^2\le\norm{a^*a}\le\norm{a^*}\norm a$, así $\norm a\le\norm{a^*}$. Sustituyendo $a$ por $a^*$ se obtiene la igualdad de normas. Por tanto $\norm{a^*a}\le\norm{a^*}\norm a=\norm a^2$, que junto con la hipótesis da la identidad $C^*$.
\end{solucion}

\begin{ejercicio}\label{ej:v3-l25-8}
Sea $(A_j)_{j\in J}$ una familia de álgebras $C^*$. Pruebe que son álgebras $C^*$, con operaciones coordenadas y norma del supremo: (i) $\prod_j^{\mathrm b}A_j=\{(a_j):\sup_j\norm{a_j}<\infty\}$; (ii) $\bigoplus_j^{c_0}A_j=\{(a_j):\text{para cada }\eps>0,\ \{j:\norm{a_j}\ge\eps\}\text{ es finito}\}$.
\end{ejercicio}
\begin{solucion}
(i) Las operaciones preservan la acotación y la norma es submultiplicativa. Una sucesión de Cauchy converge coordenada a coordenada en los espacios completos $A_j$; la cota de Cauchy uniforme pasa al límite y prueba convergencia en la norma del supremo hacia una familia acotada. La involución coordenada es isométrica y
$\norm{a^*a}=\sup_j\norm{a_j^*a_j}=\sup_j\norm{a_j}^2=\norm a^2$.

(ii) La condición dada implica acotación: fuera del conjunto finito correspondiente a $\eps=1$, las normas son menores que uno, y dentro hay un máximo finito. Sumas, productos y adjuntos conservan la condición; para productos, use $\norm{a_jb_j}\le\norm a\norm{b_j}$. Es un subespacio cerrado del producto acotado: si $a^{(n)}\to a$ uniformemente, elija $n$ con $\norm{a-a^{(n)}}<\eps/2$; entonces $\{j:\norm{a_j}\ge\eps\}\subset\{j:\norm{a_j^{(n)}}\ge\eps/2\}$, que es finito. Por tanto es completo y hereda la identidad $C^*$.
\end{solucion}

\section{Adjuntar una identidad compatible con el adjunto}\label{lec:v3-26}
\fuentelectura{26}{El cálculo funcional continuo}

Un $*$-homomorfismo es una aplicación compleja lineal, multiplicativa y que conserva la involución. Cuando se precisa preservación de la identidad se dice \emph{unital}. Más adelante veremos que entre álgebras $C^*$ la continuidad de un $*$-homomorfismo es automática.

\begin{teorema}[Unitización $C^*$]\label{v3:unitizacion-cstar}
Si $A$ es una álgebra $C^*$ no unitaria, su unitización algebraica $\widetilde A=A\oplus\C$, con
\[
 (a,\lambda)(b,\mu)=(ab+\lambda b+\mu a,\lambda\mu),\qquad
 (a,\lambda)^*=(a^*,\overline\lambda),
\]
es un álgebra $C^*$ unitaria con la norma
\[
 \norm{(a,\lambda)}_*=\sup_{\norm z\le1}\norm{az+\lambda z}.
\]
La inclusión de $A$ es isométrica y su imagen es un ideal maximal cerrado.
\end{teorema}
\begin{proof}
Las identidades de la involución se verifican coordenada a coordenada; en particular, al conjugar el producto su primera componente es $b^*a^*+\overline\lambda b^*+\overline\mu a^*$, la primera componente de $(b,\mu)^*(a,\lambda)^*$. Sea $L_{(a,\lambda)}z=az+\lambda z$. La norma propuesta es la norma de operador y $L_{xy}=L_xL_y$, de modo que es submultiplicativa.

Primero, para $a\in A$, $\norm{L_a}=\norm a$: la cota superior es inmediata, y si $a\ne0$ se prueba la inferior evaluando en $a^*/\norm a$, pues $\norm{aa^*}=\norm a^2$. Si $L_{(a,\lambda)}=0$ y $\lambda\ne0$, $-a/\lambda$ sería identidad izquierda de $A$. Una identidad izquierda $e$ en un álgebra con involución produce una identidad derecha $e^*$; como $ee^*=e^*$ por la primera propiedad y $ee^*=e$ por la segunda, ambas coinciden y $A$ sería unitaria. Por hipótesis esto es imposible. Si $\lambda=0$, $\norm{L_a}=\norm a$ da $a=0$. Así la norma es definida positiva.

Demostremos las cotas que aseguran completitud. Se tiene
\[
 |\lambda|\le\norm{L_{(a,\lambda)}},\qquad
 \norm a\le2\norm{L_{(a,\lambda)}}.
\]
Para la primera, si $\norm{L_a+\lambda I}<|\lambda|$, entonces $L_a=-\lambda[I-(L_a+\lambda I)/\lambda]$ sería invertible en $\Lcal(A)$ por Neumann. Tome $e\in A$ con $ae=a$. Para todo $b\in A$, $a(eb-b)=0$; la inyectividad de $L_a$ implica $eb=b$, contradiciendo que $A$ no tiene identidad. La segunda cota se sigue de $\norm a=\norm{L_a}\le\norm{L_a+\lambda I}+|\lambda|$. Por tanto
\[
 \norm{(a,\lambda)}_*\le\norm a+|\lambda|
 \le3\norm{(a,\lambda)}_*.
\]
La norma es equivalente a la norma suma completa, y por ello es completa.

Falta la identidad $C^*$. Para cualquier $v\in A$,
$\norm v=\sup_{\norm w\le1}\norm{wv}$, por el mismo argumento con $v^*/\norm v$. Si $x=(a,\lambda)$ y $z,w\in A$ tienen norma a lo sumo uno,
\[
 \norm{wx^*z}=\norm{z^*xw^*}\le\norm x_*.
\]
Tomar primero supremo en $w$ y luego en $z$ prueba $\norm{x^*}_*\le\norm x_*$; aplicar lo mismo a $x^*$ da igualdad. Además
\[
 \norm{xz}^2=\norm{z^*x^*xz}\le\norm{z^*}\norm{x^*xz}
 \le\norm{x^*x}_*\norm z^2.
\]
Tomar el supremo sobre $z$ da $\norm x_*^2\le\norm{x^*x}_*$. La desigualdad inversa resulta de submultiplicatividad e isometría de la involución. Finalmente, $A\oplus\{0\}$ es un ideal cerrado y el cociente, identificado con la segunda coordenada, es $\C$, luego es maximal.
\end{proof}

Para un álgebra no unitaria definimos el espectro de $a\in A$ dentro de $\widetilde A$. El cero siempre pertenece a él, porque la proyección escalar anula $a$ e impide que sea invertible. La norma original no cambia, de modo que las propiedades de elementos normales y autoadjuntos ya probadas se extienden a este caso.

\section{Representación de las álgebras conmutativas}

\begin{lema}[Permanencia espectral $C^*$]\label{v3:permanencia-cstar}
Si $B\subset A$ es una subálgebra $C^*$ cerrada que contiene la identidad de $A$, entonces $\sigma_B(b)=\sigma_A(b)$ para todo $b\in B$.
\end{lema}
\begin{proof}
Sólo hay que probar que un elemento de $B$ invertible en $A$ tiene su inversa en $B$. Primero suponga $b=b^*$ invertible en $A$. El espectro de $b$ en $A$ es un compacto real que evita cero. Elija $R>\norm b$. Por mapeo espectral polinomial y normalidad,
\[
 \norm{1-b^2/R^2}
 =\max_{t\in\sigma_A(b)}|1-t^2/R^2|<1.
\]
La serie de Neumann, cuyos términos pertenecen a $B$, muestra que $(b^2)^{-1}\in B$, y entonces $b^{-1}=b(b^2)^{-1}\in B$. Para $b$ general invertible, $b^*b$ es autoadjunto e invertible en $A$. El caso anterior da $(b^*b)^{-1}\in B$, y $b^{-1}=(b^*b)^{-1}b^*\in B$. No hemos usado positividad espectral de $b^*b$ en esta prueba, por lo que no hay circularidad con el capítulo siguiente.
\end{proof}

\begin{lema}[Stone--Weierstrass en compacto Hausdorff]\label{v3:stone-weierstrass}
Si $K$ es compacto Hausdorff y $\mathcal B\subset C(K,\C)$ es una subálgebra que contiene constantes, separa puntos y es estable bajo conjugación, entonces es uniformemente densa en $C(K,\C)$.
\end{lema}
\begin{proof}
Primero justificamos la herramienta de aproximación usada en la prueba. Para $h\in C([0,1],\R)$, los polinomios de Bernstein
\[
 B_nh(t)=\sum_{k=0}^nh(k/n)\binom nk t^k(1-t)^{n-k}
\]
convergen uniformemente a $h$. En efecto, los pesos son no negativos y suman uno; su promedio de $k/n$ es $t$ y su promedio de $(k/n-t)^2$ es $t(1-t)/n\le1/(4n)$, identidades que se obtienen diferenciando el binomio. Dado $\delta>0$, divida la suma del error entre $|k/n-t|<\delta$ y su complemento. La primera parte se acota por el módulo de continuidad $\omega_h(\delta)$; la suma de pesos de la segunda es a lo sumo $1/(4n\delta^2)$, de modo que
\[
 \norm{B_nh-h}_\infty\le\omega_h(\delta)+\frac{\norm h_\infty}{2n\delta^2}.
\]
Primero elija $\delta$ y después $n$. Por cambio afín, toda función continua real de un intervalo compacto, en particular $|t|$, admite aproximación polinomial uniforme.

Sea $\mathcal R$ la parte real de $\mathcal B$, es decir, sus elementos reales. Es una subálgebra real con constantes que separa puntos: si un elemento complejo distingue dos puntos, su parte real o su parte imaginaria también los distingue, y ambas pertenecen a $\mathcal B$ por estabilidad bajo conjugación. Su clausura uniforme $\overline{\mathcal R}$ es estable bajo valor absoluto: para $g$ en ella, aproxime $|t|$ por polinomios en un intervalo que contiene el rango de $g$, y evalúelos en $g$. Por tanto es estable bajo máximos y mínimos finitos, mediante $\max(g,h)=(g+h+|g-h|)/2$ y la fórmula análoga para el mínimo.

Fije $f\in C(K,\R)$ y $\eps>0$. Para $x,y\in K$ existe $a_{x,y}\in\mathcal R$ con $a_{x,y}(x)=f(x)$ y $a_{x,y}(y)=f(y)$: si $x\ne y$, una función que los separa se ajusta mediante una transformación afín; si coinciden, basta una constante. Para $x$ fijo, los abiertos $\{t:a_{x,y}(t)>f(t)-\eps\}$ cubren $K$ al variar $y$. Elija un subcubrimiento finito y sea $g_x$ el máximo de las correspondientes funciones. Entonces $g_x\in\overline{\mathcal R}$, $g_x>f-\eps$ en todo $K$ y $g_x(x)=f(x)$. Los abiertos $\{t:g_x(t)<f(t)+\eps\}$ cubren $K$ al variar $x$. Elija otro subcubrimiento finito y tome el mínimo $g$ de esas funciones $g_x$. Se tiene $f-\eps<g<f+\eps$ en todo $K$, y $g\in\overline{\mathcal R}$. Como esta clausura es cerrada y $\eps$ es arbitrario, $f\in\overline{\mathcal R}$. Finalmente aproxime por separado las partes real e imaginaria de una función compleja. Ningún paso utiliza una métrica en $K$.
\end{proof}

\begin{teorema}[Gelfand para álgebras $C^*$ conmutativas]\label{v3:gelfand-cstar}
Si $A$ es conmutativa y unitaria, $\mathcal G:A\to C(\Delta(A))$ es un $*$-isomorfismo isométrico sobreyectivo.
\end{teorema}
\begin{proof}
Todos los elementos son normales, por lo que $\norm{\widehat a}_\infty=r(a)=\norm a$. Así la imagen es cerrada: una sucesión de imágenes de Cauchy procede de una sucesión de Cauchy en el espacio completo $A$. Si $h=h^*$, cada carácter satisface $\varphi(h)\in\sigma(h)\subset\R$. Para $a=h+ik$ con $h,k$ autoadjuntos, se obtiene
$\varphi(a^*)=\varphi(h)-i\varphi(k)=\overline{\varphi(a)}$. Esto prueba que la imagen es estable bajo conjugación; la igualdad de conjuntos espectrales conjugados, por sí sola, no habría probado esta igualdad punto por punto.

La imagen contiene las constantes, porque $\widehat1=1$, y separa caracteres distintos, porque dos funcionales distintos difieren en algún $a$. El lema~\ref{v3:stone-weierstrass} implica que es densa, y al ser cerrada es todo $C(\Delta(A))$.
\end{proof}

\begin{proposicion}[Generadores y coordenadas]\label{v3:generador-espectro}
Si un álgebra de Banach conmutativa unitaria es la clausura de los polinomios en $a$, entonces $\varphi\mapsto\varphi(a)$ es un homeomorfismo de $\Delta(A)$ sobre $\sigma_A(a)$. Lo mismo ocurre si $A$ es $C^*$, conmutativa y unitaria, y está generada como álgebra $C^*$ unitaria por un elemento $a$.
\end{proposicion}
\begin{proof}
La aplicación es continua por la topología débil estrella y sobreyectiva por la descripción del espectro mediante caracteres. Si dos caracteres coinciden en $a$, coinciden en sus potencias, en todos los polinomios y, por continuidad, en su clausura. En el caso $C^*$ también coinciden en $a^*$ porque preservan conjugación, así que coinciden en todos los polinomios en $a,a^*$ y en su clausura. En ambos casos se obtiene una biyección continua de un compacto sobre un Hausdorff, cuya inversa es continua porque las imágenes de cerrados son compactas y por tanto cerradas.
\end{proof}

\begin{teorema}[Cálculo funcional continuo]\label{v3:calculo-continuo}
Sea $a$ normal en un álgebra $C^*$ unitaria $A$. Existe un único $*$-homomorfismo unital
\[
 C_a:C(\sigma_A(a))\longrightarrow A,\qquad f\longmapsto f(a),
\]
que envía la función coordenada $z$ a $a$. Es isométrico, su imagen es $C^*(1,a)$, extiende el cálculo holomorfo y cumple
\[
 \sigma_A(f(a))=f(\sigma_A(a)).
\]
\end{teorema}
\begin{proof}
La subálgebra $B=C^*(1,a)$, clausura de los polinomios en $a,a^*$, es conmutativa porque $a$ es normal. La permanencia espectral da $\sigma_B(a)=\sigma_A(a)$. Por la proposición anterior, $\tau:\Delta(B)\to\sigma_A(a)$, $\tau(\varphi)=\varphi(a)$, es un homeomorfismo. Defina $C_a(f)=\mathcal G_B^{-1}(f\circ\tau)$. Cada factor es un $*$-isomorfismo isométrico, por lo que se obtienen todas las propiedades algebraicas y de norma, incluida la imagen.

Si $f$ es holomorfa cerca del espectro, su integral calculada en $B$ coincide con la calculada en $A$: los resolventes pertenecen a $B$ por permanencia y son las mismas inversas. La transformada de Gelfand de esa integral es $f\circ\widehat a=f\circ\tau$, de modo que coincide con $C_a(f)$. Esta prueba evita necesitar una versión adicional del teorema de Runge. La unicidad se debe a que cualquier $*$-homomorfismo unital $F$ que envía $z$ a $a$ envía $\overline z$ a $a^*$ y coincide en los polinomios en $z,\overline z$; éstos son densos por Stone--Weierstrass. Para justificar el paso al límite sin suponer continuidad de antemano, una inversa se conserva bajo $F$, de modo que $\sigma(F(f))\subset\sigma(f)$. El elemento $F(f)$ es normal, porque $f$ conmuta con su conjugado y $F$ conserva productos y adjuntos. Por tanto $\norm{F(f)}=r(F(f))\le r(f)=\norm f_\infty$. Así $F$ es contractivo y la igualdad sobre el conjunto denso se extiende a todo el dominio.

Finalmente, la isomorfía con $C(\sigma(a))$ da el espectro de $f(a)$ dentro de $B$ como $f(\sigma(a))$. La permanencia espectral lo identifica con su espectro en $A$.
\end{proof}

\section{Ejercicios sobre funciones continuas y caracteres}

\begin{ejercicio}\label{ej:v3-l26-1}
Si $F:C(\sigma(a))\to A$ es un $*$-homomorfismo que extiende el cálculo holomorfo, pruebe $F=C_a$.
\end{ejercicio}
\begin{solucion}
La extensión implica $F(1)=1$ y $F(z)=a$; preservar el adjunto implica $F(\overline z)=a^*$. Por multiplicatividad y linealidad coincide con $C_a$ en cada polinomio en $z,\overline z$. Estos polinomios contienen las constantes, separan puntos del compacto $\sigma(a)$ y son estables bajo conjugación; Stone--Weierstrass los hace densos. La continuidad de ambas aplicaciones da la igualdad en toda función continua.
\end{solucion}

\begin{ejercicio}\label{ej:v3-l26-2}
Si $A$ es conmutativa, $C^*$ y no unitaria, pruebe que $\Delta(A)$ es localmente compacto Hausdorff y que $\mathcal G:A\to C_0(\Delta(A))$ es un $*$-isomorfismo isométrico.
\end{ejercicio}
\begin{solucion}
Sea $K=\Delta(\widetilde A)$. La proyección escalar $\chi_\infty(a,\lambda)=\lambda$ es un carácter. Cada carácter no nulo $\varphi$ de $A$ se extiende de manera única a $\widetilde\varphi(a,\lambda)=\varphi(a)+\lambda$. Recíprocamente, un carácter de $\widetilde A$ distinto de $\chi_\infty$ se restringe a un carácter no nulo de $A$. Las fórmulas de evaluación prueban que esta biyección identifica la topología débil estrella de $\Delta(A)$ con la de $K\setminus\{\chi_\infty\}$, que es localmente compacta Hausdorff.

La representación unitaria $\widetilde A\simeq C(K)$ lleva $A=\ker\chi_\infty$ al ideal de funciones que se anulan en $\chi_\infty$. Restringirlas a $K\setminus\{\chi_\infty\}$ identifica ese ideal con $C_0$ de dicho espacio: si $f(\chi_\infty)=0$, los conjuntos $\{|f|\ge\eps\}$ son compactos y evitan ese punto; recíprocamente, una función que se anula en infinito se extiende continuamente por cero, porque fuera de un compacto su módulo es menor que cualquier $\eps$. Todos estos mapas preservan productos, adjuntos y normas. El punto $\chi_\infty$ no es aislado: si lo fuera, su complemento sería compacto abierto y el indicador de ese complemento correspondería a una identidad de $A$, contradicción.
\end{solucion}

\begin{ejercicio}\label{ej:v3-l26-3}
Para un espacio compacto Hausdorff $X$, pruebe que $X$ es totalmente disconexo si y sólo si el subespacio generado por las proyecciones de $C(X)$ es denso.
\end{ejercicio}
\begin{solucion}
Una proyección de $C(X)$ toma sólo los valores $0,1$, por lo que es el indicador de un conjunto abierto y cerrado. Recordemos por qué en un compacto Hausdorff totalmente disconexo estos conjuntos separan puntos. Para $x\in X$, sea $Q_x$ la intersección de sus vecindades abiertas y cerradas. Afirmamos que $Q_x$ es conexo. Si se separara en dos compactos no vacíos $P,Q$, elija abiertos disjuntos $U,V$ que los contengan. Cada punto del compacto $X\setminus(U\cup V)$ es excluido por alguna vecindad abierta y cerrada de $x$; una intersección finita de ellas da un conjunto $C$ abierto y cerrado con $Q_x\subset C\subset U\cup V$. Los conjuntos $C\cap U$ y $C\cap V$ son abiertos y cerrados en $X$, y el que contiene a $x$ excluiría una parte de $Q_x$, contradicción. Así $Q_x$ es conexo y, por total desconexión, $Q_x=\{x\}$. Para $y\ne x$ existe entonces un conjunto abierto y cerrado que contiene a $x$ y no a $y$.

Los indicadores de esos conjuntos separan puntos. Su espacio generado es un álgebra (el producto de dos indicadores es el indicador de la intersección), contiene las constantes y es estable bajo conjugación; Stone--Weierstrass da densidad. Recíprocamente, una proyección es constante en cualquier subconjunto conexo, porque su imagen continua está en $\{0,1\}$. Si dos puntos distintos pertenecieran a un mismo subconjunto conexo, todas las combinaciones de proyecciones y sus límites uniformes tendrían iguales valores allí. Esto contradice que las funciones continuas de un compacto Hausdorff separan puntos. Por tanto toda componente conexa es un singleton.
\end{solucion}

\begin{ejercicio}\label{ej:v3-l26-4}
Si $A$ es una álgebra $C^*$ posiblemente no conmutativa y $a\in A$ es normal, construya su cálculo funcional continuo que extiende al holomorfo.
\end{ejercicio}
\begin{solucion}
En el caso unitario la construcción completa está en el teorema~\ref{v3:calculo-continuo}: se trabaja en $C^*(1,a)$, que es conmutativa, y la permanencia espectral asegura que se usa el espectro correcto. Si $A$ no tiene identidad, se realiza la construcción en $\widetilde A$. Para $f(0)=0$, el elemento $f(a)$ pertenece a $A$: la proyección escalar aplicada al cálculo continuo vale $f(0)$, primero para polinomios en $z,\overline z$ y luego por continuidad. Así el cálculo con valores en $A$ tiene como dominio natural $\{f\in C(\sigma_{\widetilde A}(a)):f(0)=0\}$. Las funciones con $f(0)\ne0$ dan en general elementos de la unitización. Esta precisión es necesaria si el enunciado no supone identidad.
\end{solucion}

\begin{ejercicio}\label{ej:v3-l26-5}
Sea $A$ una álgebra $C^*$ conmutativa unitaria generada por $a_1,\ldots,a_n$ (algebraicamente y luego por clausura, o como álgebra $C^*$). Pruebe que existe un compacto $X\subset\C^n$ y un $*$-isomorfismo $A\to C(X)$ que envía $a_k$ a la coordenada $z_k$.
\end{ejercicio}
\begin{solucion}
Defina $\tau:\Delta(A)\to\C^n$ por $\tau(\varphi)=(\varphi(a_1),\ldots,\varphi(a_n))$ y sea $X=\tau(\Delta(A))$. La continuidad de las coordenadas hace a $X$ compacto. Si dos caracteres coinciden en los generadores, coinciden en todos sus polinomios, en los adjuntos cuando éstos se incluyen, y por continuidad en todo $A$. Por ello $\tau$ es inyectiva y es un homeomorfismo sobre $X$. El isomorfismo buscado es $a\mapsto\widehat a\circ\tau^{-1}$. Evaluarlo en un generador da exactamente $z_k$. La hipótesis $C^*$ procede del contexto de la lista y no puede suprimirse: una álgebra de Banach conmutativa general puede tener radical no nulo y no ser isomorfa a un espacio $C(X)$.
\end{solucion}

\begin{ejercicio}\label{ej:v3-l26-6}
Sea $X$ compacto Hausdorff, $x_0\in X$, y sea $A$ el conjunto de sucesiones $(f_n)$ acotadas en $C(X)$ tales que $(f_n(x_0))$ converge. Pruebe que $A$ es una álgebra conmutativa unitaria y determine su espacio de caracteres.
\end{ejercicio}
\begin{solucion}
Con la norma $\norm f=\sup_n\norm{f_n}_\infty$, el producto acotado $B=\prod_{n\ge1}^{\mathrm b}C(X)$ es un álgebra $C^*$ conmutativa unitaria. La aplicación $q:B\to\ell^\infty$, $q(f)=(f_n(x_0))$, es un $*$-homomorfismo contractivo. Como el espacio $c$ de sucesiones convergentes es cerrado en $\ell^\infty$, $A=q^{-1}(c)$ es una subálgebra $C^*$ cerrada y contiene a la identidad. La cerradura de $c$ se verifica directamente: una aproximación uniforme por una sucesión convergente hace que la sucesión límite sea de Cauchy en sus índices.

Para describir todos los caracteres, incluidos los que no son evaluaciones en una coordenada finita, sea $K=\Delta(B)$ y defina $\gamma_n\in K$ por $\gamma_n(f)=f_n(x_0)$. Considere el compacto no vacío
\[
 F=\bigcap_{N\ge1}\overline{\{\gamma_n:n\ge N\}}^{\,K}.
\]
Es no vacío por la propiedad de intersección finita en el compacto $K$. Afirmamos que, bajo $B\simeq C(K)$, la subálgebra $A$ corresponde exactamente a las funciones continuas constantes sobre $F$. Si $f_n(x_0)\to L$, todo punto de adherencia de las colas tiene valor $L$ por continuidad. Recíprocamente, si $\widehat f=L$ en $F$ y la sucesión no convergiera a $L$, habría $\eps>0$ e infinitos índices con $|\gamma_n(f)-L|\ge\eps$. Las clausuras de las colas de esa subsucesión tendrían un punto común en $F$ donde, por cerradura de la desigualdad, el valor se alejaría de $L$ al menos $\eps$, contradicción.

Sea $K/F$ el espacio que identifica todos los puntos de $F$ en un único punto y deja distintos los demás. Es compacto Hausdorff: la relación de equivalencia es el conjunto cerrado diagonal unido a $F\times F$; alternativamente, la normalidad de $K$ separa el cerrado $F$ de cualquier punto exterior y permite construir vecindades disjuntas en el cociente. Las funciones continuas constantes en $F$ son exactamente las funciones que descienden a $K/F$. En consecuencia
\[
 \boxed{\ \Delta(A)\simeq K/F,\qquad K=\Delta\!\left(\prod_{n\ge1}^{\mathrm b}C(X)\right).\ }
\]
Esta descripción es concreta en términos de la compactificación de Stone--Čech: $B=C_b(\N\times X)$, con $\N$ discreto, así $K=\beta(\N\times X)$, y $F$ es el resto en infinito de la copia $\N\times\{x_0\}$. Por tanto se colapsa \emph{sólo ese resto} a un punto. El carácter del punto colapsado es $f\mapsto\lim_n f_n(x_0)$; cada $(n,x)$ da el carácter $f\mapsto f_n(x)$, y pueden existir otros caracteres en el resto no colapsado.

Para justificar la última advertencia, si $x_1\ne x_0$, elija $g\in C(X)$ con $g(x_0)=0$, $g(x_1)=1$. Toda sucesión escalar acotada $(c_n)$ produce $(c_ng)\in A$. Para un ultrafiltro libre $\mathcal U$ sobre $\N$, la fórmula $f\mapsto\lim_{n\to\mathcal U}f_n(x_1)$ es un carácter, porque los límites de ultrafiltro de sucesiones acotadas respetan sumas y productos. Distintos ultrafiltros se distinguen con sucesiones indicadoras y, por tanto, con elementos $(c_ng)$ de $A$. En cambio, si $X=\{x_0\}$, $A=c$ y el cociente se reduce a $\N\cup\{\infty\}$, la compactificación por un punto.
\end{solucion}

La norma ya permite evaluar cualquier función continua del espectro de un elemento normal. En particular, funciones reales no negativas y raíces cuadradas producirán una estructura de orden. Ese paso será la base para estudiar funcionales positivos y construir representaciones por operadores. Como apoyo adicional véanse \cite{maxcstar,maxcalculo,murphy1990,zhu1993}.
